\documentclass[11pt]{article}
\usepackage[T1]{fontenc}
\usepackage{lmodern}
\usepackage[margin=0.85in]{geometry}
\usepackage{amsmath,amssymb,mathtools,booktabs,array,microtype,xcolor,xurl,hyperref,fancyhdr}
\hypersetup{colorlinks=true,linkcolor=blue!45!black,urlcolor=blue!45!black,
 pdftitle={Finite value of the submitted gluon bubble},pdfauthor={Yang--Mills group}}
\setlength{\parindent}{0pt}\setlength{\parskip}{5pt plus 1pt}
\setlength{\headheight}{14pt}
\pagestyle{fancy}\fancyhf{}\fancyhead[L]{\small Submitted gluon bubble}
\fancyhead[R]{\small Finite identity verified in Lean}\fancyfoot[C]{\thepage}
\newcommand{\R}{\mathbb R}
\newcommand{\code}[1]{\ifmmode\text{\small\nolinkurl{#1}}\else{\small\nolinkurl{#1}}\fi}
\newcommand{\E}{\mathbb E}
\newcommand{\dd}{\,\mathrm d}
\begin{document}
\begin{center}
{\LARGE\bfseries Finite value of the submitted gluon bubble}\par
\vspace{4pt}
November 4 Yang--Mills submission\par
25 November 2026
\end{center}

This note derives the finite value of the actual gauge/gauge bubble,
\texttt{TwoPoint.bubble G false false}, including its dependence on the retained
infrared cutoff. It uses the submitted heat kernels, dimensional continuation
and Taylor subtraction. The resulting infrared limit agrees with the earlier
conventional integral calculation.

\textbf{Verification status.} The full infrared finite-value identity is
now proved in Lean. The theorem \code{GluonSolution.acceptance} discharges
\code{GluonLoop.SubmittedFinitePartTest} for the actual submitted evaluator,
for every nonzero physical momentum, tensor component and color, and every
positive $\bar\mu$, under the submitted simple-group hypothesis. Its proof
includes the extra polarization trace, radial continuation, regulator
coefficient and ultraviolet subtraction. The retained-cutoff equality is
also proved for the two unit-momentum probes. The proof uses only
\code{propext}, \code{Classical.choice} and \code{Quot.sound}; it assumes
no finite-part identity. Human review remains minimal. These checks concern
this graph contribution, within the submitted definitions, and form part
of the continuing review.

\section{The identity}

Strip off the overall coupling $g^2$. Let $Q=p^2>0$,
$C^{ab}=\delta^{ab}C_A$, and
\[
 P^{\rm T}_{ij}=\delta_{ij}-\frac{p_i p_j}{Q},\qquad
 P^{\rm L}_{ij}=\frac{p_i p_j}{Q},\qquad
 L=\log\frac{\bar\mu^2}{Q},\qquad X=Qe^{2s}.
\]
The conventions are Euclidean Landau gauge and the submitted \emph{log-density}
sign. The induced effective-action kernel has the opposite sign.
Set $y=\log\bar\mu+(\gamma-\log(4\pi))/2$, and write
\begin{equation}
 \begin{aligned}
 \mathcal B_G&=(\code{TwoPoint.bubble}\ G\ \code{false}\ \code{false}).
 \code{instructions.renorm},\\
 \mathcal F^{ab}_{ij}(s,\bar\mu;p)&=
 [\code{Outlook.four}\ 4\,\mathcal B_G\ (y,s)]
 (\code{TwoPoint.word}\ i\ j\ a\ b\ p).
 \end{aligned}
 \label{eq:target}
\end{equation}
For the submitted simple-group hypothesis, the calculation gives
\begin{equation}
 \boxed{\displaystyle
 \mathcal F^{ab}_{ij}
 =\frac{C^{ab}Q}{16\pi^2}
 \left[
 \left(\mathcal P_{\rm T}(X)+\frac{25}{12}L+\frac{89}{36}\right)P^{\rm T}_{ij}
 +\left(\mathcal P_{\rm L}(X)-\frac14 L-\frac12\right)P^{\rm L}_{ij}
 \right].}
 \label{eq:cutoff}
\end{equation}
Here $E_1(x)=\int_x^\infty e^{-t}\dd t/t$, and the cutoff terms are
\begin{align}
 \mathcal P_{\rm T}(X)
 &=\frac{14}{3X^3}-\frac{25}{2X^2}-\frac6X
 +e^{-X/2}\left(-\frac{14}{3X^3}+\frac{61}{6X^2}+\frac{181}{24X}\right)
 -\frac{25}{12}E_1(X/2),\label{eq:PT}\\
 \mathcal P_{\rm L}(X)
 &=-\frac{14}{X^3}+\frac{9}{2X^2}-\frac{9}{2X}
 +e^{-X/2}\left(\frac{14}{X^3}+\frac{5}{2X^2}-\frac{1}{8X}\right)
 +\frac14 E_1(X/2).\label{eq:PL}
\end{align}
Both tend to zero as $X\to\infty$. Thus, with dimensional subtraction taken
first as in \eqref{eq:target}, the full finite-value limit is
\begin{equation}
 \boxed{\displaystyle
 \lim_{s\to\infty}\mathcal F^{ab}_{ij}
 =\frac{C^{ab}}{16\pi^2}
 \left[
 \delta_{ij}Q\left(\frac{25}{12}L+\frac{89}{36}\right)
 -p_i p_j\left(\frac73 L+\frac{107}{36}\right)
 \right].}
 \label{eq:limit}
\end{equation}
This concerns the gluon bubble alone. It includes neither the ghost bubble
nor a tadpole. No cancellation of a retained-cutoff tadpole is assumed.

\section{What is being contracted}

The source cubic vertex is evaluated in
\code{YM/WilsonMS/Beta/Flash.lean}, \code{field_W}.
The review extends that evaluation to every auxiliary sector.
In ordinary tensor notation the vertex, with its common
color and imaginary factors removed, is
\[
 V_{e,\alpha\beta}
 =e_\alpha(p-k)_\beta+(2e\cdot k+e\cdot p)\delta_{\alpha\beta}
 -(k+2p)_\alpha e_\beta .
\]
The source uses $\operatorname{precision}_{\alpha\beta}(k)
=k^2\delta_{\alpha\beta}-k_\alpha k_\beta$ and the shell weight
$t^3e^{-t^2k^2}$. Put $A=k^2$, $B=(k+p)^2$. The heated numerator is the
contraction of two copies of $V$ with $\operatorname{precision}(k)$ and
$\operatorname{precision}(k+p)$, so it is $AB$ times the corresponding
projector numerator.

Use a unit vector $e$ and first take $p=x e_1$. For a transverse probe
$e=e_2$, put $K=k^2-k_1^2$; for the longitudinal probe use $e=e_1$.
The finite index sums reduce to
\begin{align}
 N_{\rm T}&=4\{x^2(A+B-k_2^2)K+(H-1)ABk_2^2\},\label{eq:NT}\\
 N_{\rm L}&=(2k_1+x)^2\{(H-1)AB-x^2K\}.\label{eq:NL}
\end{align}
$H$ denotes the vector-trace dimension, kept separate from the Gaussian
integration dimension until the continuation is performed. The source
sewing supplies a factor $2$ and two orientations. After the even momentum
integral their combined factor is $4C^{ab}$. This normalization follows
from \code{gauge_hit}, \code{exchangeW} and \code{gauge_half}; the review's
\code{GluonLoopTrace.trace_value} proves their combined trace formula in
arbitrary sectors. It is not an inserted conventional symmetry factor.

\subsection{The extra polarization}

If both momenta and the external polarization have zero component along a
new axis, that axis contributes
\begin{equation}
 AB(2e\cdot k+e\cdot p)^2
 \label{eq:spectator}
\end{equation}
to the heated numerator. Indeed, the two vertex indices must both equal
the new axis; the surviving vertex entry is $2e\cdot k+e\cdot p$, and the
two precision entries are $A$ and $B$. This is precisely the extra-axis
contraction called \code{yeast}.

The physical external sector has four axes. Inside \code{loopCarpet}, the
extra radius is placed in sector 1, which has six test axes. Consequently
the source prescription is
\[
 \code{innerTrace}=\text{six-axis trace}+2(z-1)\,\code{yeast},
 \qquad H=6+2(z-1)=4+2z.
\]
Replacing the coefficient by $2z$ at this stage would count two extra
polarizations. The remaining closing operation is the source's radial
integral and radial derivative; agreement at integer dimensions alone is
not being used to infer its continuation.

\section{Evaluation of the submitted heat integral}

Take a unit reference momentum $e_1$ and vary the external ray $p=xe_1$.
Introduce
\[
 u=\frac{t^2}{1+t^2}\in[0,\tfrac12],\quad
 a=-u,\quad b=1-u,\quad c=a+b=1-2u,\quad
 \rho=\frac{1-u}{2},\quad X=x^2.
\]
Completing the square shifts $k$ to $\ell+a x e_1$. The normalized
Gaussian has covariance $\rho\delta_{ij}$; its integration dimension is
denoted by $D$. For compactness define
\begin{align*}
 U_0&=D(D+2)\rho^2+[D(a^2+b^2)+4ab]X\rho+a^2b^2X^2,\\
 U_{\rm T}&=(D+2)(D+4)\rho^3+
 [(D+2)(a^2+b^2)+4ab]X\rho^2+a^2b^2X^2\rho,\\
 U_{\rm L}&=(D+2)(D+4)\rho^3+
 [(D+2)(a^2+b^2)+12ab]X\rho^2+a^2b^2X^2\rho.
\end{align*}
These are respectively $\E[AB]$, $\E[AB\ell_2^2]$ and
$\E[AB\ell_1^2]$. Wick contraction of \eqref{eq:NT}--\eqref{eq:NL} gives
\begin{align}
 M_{\rm T}(D,H)&=4\{X[(2D^2+D-5)\rho^2+
 (a^2+b^2)X(D-1)\rho]+(H-1)U_{\rm T}\},\label{eq:MT}\\
 M_{\rm L}(D,H)&=(H-1)\{4U_{\rm L}+8c^2(D+2)X\rho^2
 +8abc^2X^2\rho+c^2XU_0\}\notag\\[-2pt]
 &\hspace{1cm}-X\{4(D-1)\rho^2+c^2X(D-1)\rho\}.\label{eq:ML}
\end{align}

\subsection{Why these moments give the source continuation}

In split coordinates the physical Gaussian integral gives
$M_r(4,H)+R^2B_r+R^4C_r$, where $r\in\{{\rm T},{\rm L}\}$ and
\begin{align*}
 B_{\rm T}&=4\{X[13\rho+(a^2+b^2)X]+(H-1)[12\rho^2+(a^2+b^2)X\rho]\},\\
 C_{\rm T}&=4[2X+(H-1)\rho],\\
 B_{\rm L}&=(H-1)\{48\rho^2+[16c^2+4(a^2+b^2)]X\rho+
 c^2(a^2+b^2)X^2\}-X(4\rho+c^2X),\\
 C_{\rm L}&=(H-1)(4\rho+c^2X).
\end{align*}
With $\nu=1+t^2$, the normalized source radial operation is
\begin{multline}
 -\frac1{\Gamma(z+1)}\int_0^\infty R^{2z+1}
 \frac1R\frac{\mathrm d}{\mathrm dR}
 \left[e^{-\nu R^2}(A_0+B_0R^2+C_0R^4)\right]\dd R\\
 =\nu^{-z}\left[A_0+\frac{z}{\nu}B_0+
 \frac{z(z+1)}{\nu^2}C_0\right],\qquad \Re z>-1.
 \label{eq:radial}
\end{multline}
Since $\rho=(2\nu)^{-1}$, the bracket equals
$M_r(4,H)+2z\rho B_r+4z(z+1)\rho^2C_r=M_r(4+2z,H)$.
The last two polynomial identities are proved in
\code{GluonHeat.split_radial_T} and \code{split_radial_L}; the integral
identity \eqref{eq:radial} uses the same Gaussian/Gamma calculation as the
already verified ghost continuation. The extra trace then sets $H=4+2z$.

The measure factor from \code{tenorK} is $(4\pi)^{-z}$ at $y=0$.
Also $t^3(1+t^2)^{-2}\dd t=u\dd u/[2(1-u)]$.
Including the factor $4$ from sewing, the coefficient of
$C^{ab}/(16\pi^2)$ in the source force is therefore
\begin{equation}
 w_r(z,X)=(4\pi)^{-z}\int_0^{1/2}(1-u)^z e^{-uX}
 \frac{2u}{1-u}M_r(4+2z,4+2z)\dd u.
 \label{eq:continued}
\end{equation}
This derivation evaluates the specified radial continuation, rather than
substituting $D$ into a four-dimensional answer. All exchanges of the
Gaussian and compact heat integrals, and the derivatives at $z=0$ and
$x=0$, are justified by polynomial Gaussian bounds. Here $1\le\nu\le2$;
the radial integrand is locally bounded by a constant times
$R^{2\Re z+1}$ at zero and decays exponentially at infinity. These bounds
are uniform for $z$ in a compact neighborhood of zero and bounded $x$.

\subsection{The force at zero regulator}

Let $h_r(u,X)=2uM_r(4,4)/(1-u)$. Directly from the moments,
\begin{align*}
 h_{\rm T}&=144u(1-u)^2+(98u-194u^2+192u^3-96u^4)X\\
 &\hspace{1cm}+(12u-24u^2+36u^3-24u^4+12u^5)X^2,\\
 h_{\rm L}&=144u(1-u)^2+(138u-714u^2+1152u^3-576u^4)X\\
 &\hspace{1cm}+(9u-96u^2+348u^3-504u^4+252u^5)X^2\\
 &\hspace{1cm}+(6u^3-30u^4+48u^5-24u^6)X^3.
\end{align*}
Thus $f_r(X)=w_r(0,X)=\int_0^{1/2}e^{-uX}h_r(u,X)\dd u$ is
\begin{align}
 f_{\rm T}(X)&=-\frac{28}{X^2}+\frac{50}{X}+12
 +e^{-X/2}\left(\frac{28}{X^2}-\frac{36}{X}-\frac{101}{4}-\frac{27X}{8}\right),\\
 f_{\rm L}(X)&=\frac{84}{X^2}-\frac{18}{X}+9
 +e^{-X/2}\left(-\frac{84}{X^2}-\frac{24}{X}-\frac94-\frac{3X}{8}\right).
\end{align}
The apparent singularities at zero are removable. Both constant terms
are $33/4$; their linear coefficients are $25/6$ and $-1/2$.
The elementary integral equality for both channels, and any linear
combination of them, is proved in \code{GluonHeat.heatPolynomial_integral}.

\section{The regulator derivative fixes the finite subtraction}

The source \code{highTensor} is minus one half of the quadratic external
momentum Taylor coefficient of the force. At unit momentum, after stripping
the common factor $C^{ab}/(16\pi^2)$, its values at zero are
\[
 H_{\rm T}(0)=-\frac{25}{12},\qquad H_{\rm L}(0)=\frac14.
\]
To compute $J_r=H'_r(0)$, let $j_r(u)$ be the coefficient of $X$ in
$e^{-uX}h_r(u,X)$, and let $\dot j_r(u)$ be the derivative of this
coefficient through $D=H=4+2z$, before differentiating $(1-u)^z$.
The polynomials are
\begin{align*}
 j_{\rm T}&=98u-338u^2+480u^3-240u^4,&
 \dot j_{\rm T}&=104u-372u^2+536u^3-268u^4,\\
 j_{\rm L}&=138u-858u^2+1440u^3-720u^4,&
 \dot j_{\rm L}&=158u-962u^2+1608u^3-804u^4.
\end{align*}
Their integrals are
\[
\setlength{\arraycolsep}{12pt}
\begin{array}{c|cc}
 &\displaystyle\int_0^{1/2}j_r\dd u
 &\displaystyle\int_0^{1/2}[\dot j_r+j_r\log(1-u)]\dd u\\[3pt]\hline
 {\rm T}&25/6&-5/36+(25/6)\log2\\[3pt]
 {\rm L}&-1/2&5/6-(1/2)\log2
\end{array}
\]
In particular, differentiating the measure factor in \eqref{eq:continued}
gives
\begin{align}
 J_{\rm T}&=\frac5{72}-\frac{25}{12}\log2+\frac{25}{12}\log(4\pi),\\
 J_{\rm L}&=-\frac5{12}+\frac14\log2-\frac14\log(4\pi).
 \label{eq:jet}
\end{align}
These constants are derived from the regulated tensor; they are not chosen
to produce the conventional answer. The source's holomorphic coefficient
\texttt{ruby ... 1} equals this real-axis derivative, by the regularity of
the continued force near zero.

For reproducibility, polynomial/logarithmic integration uses only
\[
 \int v^n\log v\dd v=
 v^{n+1}\left(\frac{\log v}{n+1}-\frac1{(n+1)^2}\right)
\]
after $v=1-u$. The exact rational script checks the displayed polynomials
and constants. The source identification is now proved in
\code{GluonJetParameter.submitted_jet_unit}; its logarithmic integral is
proved by differentiating an explicit primitive on $[0,1/2]$.

\section{Fixing the integration constant and removing the cutoff}

The source identity \code{twoVertex_laurent} identifies
\eqref{eq:target} with \code{oneLoopFinite}. The source scale derivative
is proved for this gluon force in \code{GluonLoop.submitted_finite_derivative}.
Combining that result with the source-connected Lean force evaluation gives
\[
 \frac{\mathrm d F_r(s,y)}{\mathrm ds}=e^{-2s}f_r(e^{2s})
 \quad(Q=1),
\]
where $F_r$ denotes the coefficient of $C^{ab}/(16\pi^2)$. The functions
\eqref{eq:PT}--\eqref{eq:PL} satisfy
\[
 \mathcal P'_r(X)=\frac{f_r(X)}{2X^2},\qquad
 \lim_{X\to\infty}\mathcal P_r(X)=0.
\]
Their behavior at zero is
\begin{align*}
 \mathcal P_{\rm T}(X)&=-\frac{33}{8X}+\frac{25}{12}\log X+c_{\rm T}+o(1),
 &c_{\rm T}&=-\frac{173}{72}-\frac{25}{12}(\log2-\gamma),\\
 \mathcal P_{\rm L}(X)&=-\frac{33}{8X}-\frac14\log X+c_{\rm L}+o(1),
 &c_{\rm L}&=\frac1{12}+\frac14(\log2-\gamma).
\end{align*}
The derivative and both endpoint statements are proved in Lean. The
logarithmic constant uses $E_1(x)=-\gamma-\log x+o(1)$, with the Euler
constant retained.

The source's lower Taylor subtraction gives
$-33/(8X)$ in the ultraviolet boundary of the ordinary part.
There is no linear external-momentum term, since the force is even in $x$.
The remaining, degree-at-least-three Taylor remainder tends to zero after
rescaling: a smooth $O(x^3)$ remainder becomes $O(x)$ when multiplied by
$x^{-2}$. This is the boundary encoded by \code{ordinary} and
\code{frame_lim}. The resonant quadratic term is handled separately in
\code{oneLoopFinite}. Consequently
\[
 F_r(s,y)=\mathcal P_r(e^{2s})-c_r+J_r-2yH_r(0).
\]
Indeed, the two sides have the same derivative in $s$; their difference
has zero ultraviolet limit, so their additive constants agree. Inserting
the values above gives
\begin{align*}
 -c_{\rm T}+J_{\rm T}-2yH_{\rm T}(0)
 &=\frac{25}{12}[2y-\gamma+\log(4\pi)]+\frac{89}{36},\\
 -c_{\rm L}+J_{\rm L}-2yH_{\rm L}(0)
 &=-\frac14[2y-\gamma+\log(4\pi)]-\frac12.
\end{align*}
The $\log2$ terms cancel. This cancellation and the finite rational
constants are checked in the Lean lemma \code{finite_constants}.

The source rotation law reconstructs the two tensor projectors from these
two probes. Its scale law replaces $X$ by $Qe^{2s}$, $y$ by
$y-\tfrac12\log Q$, and multiplies the unit-momentum tensor by $Q$.
The review checks both covariance and this scale law for the actual gluon
finite reading in \code{GluonLoopTensor.lean}. Setting the stated value of
$y$ proves the analytic identity \eqref{eq:cutoff}, and the large-$X$
limits yield \eqref{eq:limit}. The full limit, including covariance and
rescaling, is the conclusion of \code{GluonSolution.acceptance}. No interchange of an infrared limit with
a Laurent expansion is needed.

\section{The vector-trace check and formal verification}

Freezing only $H=4$ leaves the physical force $f_r$ and its primitive
unchanged. The error appears in the regulator derivative. The contribution
of $2\partial_H$ to the integrated quadratic force coefficient is $-2/3$
in the transverse channel and zero in the longitudinal channel. Since
$H_r$ is minus one half of that coefficient, keeping the continued vector
trace adds
\[
 \frac{C^{ab}}{16\pi^2}\,\frac13
 (Q\delta_{ij}-p_i p_j)
\]
to the finite value. Freezing it gives $77/36$ instead of $89/36$ in the
gluon transverse constant. The pole is unchanged. Adding the previously
verified ghost bubble gives the transverse finite coefficient
$(13/6)L+97/36$; the longitudinal terms cancel. These statements do not
evaluate a source tadpole, and the frozen trace is only a diagnostic
replacement, not a specification of an alternative scheme.

The proof follows the submitted evaluator through the following stages:
\begin{center}\small
\begin{tabular}{>{\raggedright\arraybackslash}p{0.35\linewidth}p{0.57\linewidth}}
\toprule
Review modules & What is proved\\\midrule
\mbox{\code{Vertex}}, \mbox{\code{Trace}}, \mbox{\code{Contraction}}, \mbox{\code{Yeast}} & Actual cubic sewing, six-axis contractions and the extra-axis trace.\\[3pt]
\mbox{\code{Gaussian}}, \mbox{\code{Carpet}}, \mbox{\code{Radial}}, \mbox{\code{Continuation}} & Physical Gaussian integration and the submitted radial continuation, on a real neighborhood of $z=0$.\\[3pt]
\mbox{\code{Jet}}, \mbox{\code{JetParameter}}, \mbox{\code{Infrared}} & The actual regulator coefficient and the explicit physical force.\\[3pt]
\mbox{\code{Flow}}, \mbox{\code{Unit}}, \mbox{\code{Tensor}}, \mbox{\code{Solution}} & Ultraviolet boundary, exact unit-probe cutoff value, covariance, rescaling and the full finite limit.\\[3pt]
\mbox{\code{Finite}}, \mbox{\code{Heat}} & Fixed target, primitive and endpoints, polynomial reduction and compact heat integral.\\
\bottomrule
\end{tabular}
\end{center}
Every module name in the table has the prefix \code{GluonLoop} and
extension \code{.lean}. The full acceptance theorem uses the same target
definition as the earlier component checks. A fresh 33-module review build
and a frozen-target Comparator run check it with the standard Lean axioms
and no permitted definition holes. The received libraries are reused;
this is not a new build of the entire submission. All 37,673 received
source files are checked against the pinned manifest before and after
Comparator. The run uses upstream's macOS development wrapper and the
Lean default kernel; it is not a sandboxed run or a second-kernel check.

Reproduce the review build with \code{scripts/build_gluon_finite.py}
and the frozen-target comparison with
\code{scripts/run_gluon_finite_comparator.py}. Their receipts are
\code{evidence/matching/gluon-finite-build.json} and the files in
\code{evidence/matching/gluon-finite-comparator/}.
The earlier 13 component audits are retained as historical evidence.
The exact-rational heat calculation remains an independent cross-check;
it is not a premise of the Lean theorem.

\smallskip
\textbf{Source provenance.} This calculation uses the November source
revision \code{b5e288edeac43ac088a2a291ac3841bdda9f20b8}.
The main source anchors are
\code{YM/MS/ScaleForest/Glue.lean} (\code{innerTrace}, \code{loopCarpet},
\code{closeG}); \code{TracedCubic.lean} and \code{LoadB.lean}
(actual gauge sewing); \code{Bubbles.lean} (\code{ratePair}, \code{bubble});
and \code{PageOne.lean} (\code{ordinary}, \code{highTensor},
\code{oneLoopFinite}). The source files are unchanged. This is a review of
the specified perturbative evaluator, not an assessment of the complete
Yang--Mills or QCD existence argument.
\end{document}
